\section*{الفصل \ref{ch:g11:lithosphere-asthenosphere} --- الغلاف الصخري والغلاف اللدن}

\begin{solution}{exo:g11:lithosphere-asthenosphere:1}
\omterm{def:g11:lithosphere-asthenosphere:lithosphere}{الغلاف الصخري} هو الغلاف الخارجي الصلب للأرض، ويتكوّن من القشرة
وأعلى الوشاح؛ \omterm{def:g11:lithosphere-asthenosphere:lithosphere}{والغلاف اللدن} هو الوشاح الأسخن والأضعف الذي
تحته، وهو ينساب ببطء على مدى أزمنة طويلة. والقشرة المحيطية والقشرة
القارية وأعلى الوشاح تنتمي إلى \omterm{def:g11:lithosphere-asthenosphere:lithosphere}{الغلاف الصخري}؛ أما اللب الخارجي
فلا.
\end{solution}

\begin{solution}{exo:g11:lithosphere-asthenosphere:2}
يغوص الجبل الجليدي حتى يزيح من الماء ما يساوي وزنه: فالكسر
الواقع تحت الماء $910/1025 \approx 0.888$، أي نحو \qty{89}{\percent}.
والكسر الذي فوقه نحو \qty{11.2}{\percent}، فيكون الارتفاع الكلي
$20/0.112 \approx \qty{180}{m}$، يقع منها نحو \qty{160}{m} تحت
الماء.
\end{solution}

\begin{solution}{exo:g11:lithosphere-asthenosphere:3}
$15 + 26 \times 3 = \qty{93}{\celsius}$. وعند \qty{100}{km} كان التدرج نفسه
سيعطي أكثر من \qty{2600}{\celsius}، أي أعلى بكثير من درجة انصهار
صخر الوشاح؛ لكن \omterm{def:g11:seismic-structure:body-waves}{الموجات S} تعبر الوشاح العلوي، فهو إذن صلب.
فلا بد أن يصغر التدرج كثيرًا مع العمق.
\end{solution}

\begin{solution}{exo:g11:lithosphere-asthenosphere:4}
النطاق من 0 إلى \qty{500}{m} فوق مستوى سطح البحر، وفيه نحو
\qty{16}{\percent} من السطح، والنطاق من 4000 إلى \qty{4500}{m}
تحت مستوى سطح البحر، وفيه نحو \qty{11}{\percent} (والنطاق من 4500 إلى
\qty{5000}{m} تحته يكاد يكون بالقدر نفسه).
\end{solution}

\begin{solution}{exo:g11:lithosphere-asthenosphere:5}
عند $P = 0$: \qty{1121}{\celsius}. وعند \qty{3}{GPa}: $-5.104 \times 9 +
132.899 \times 3 + 1120.661 \approx -45.9 + 398.7 + 1120.7 \approx
\qty{1473}{\celsius}$.
\end{solution}

\begin{solution}{exo:g11:lithosphere-asthenosphere:6}
نحو \qty{730}{\celsius} عند \qty{50}{km} (\omterm{def:g11:lithosphere-asthenosphere:lithosphere}{الغلاف الصخري})، ونحو
\qty{1370}{\celsius} عند \qty{150}{km} و\qty{1410}{\celsius} عند
\qty{250}{km} (\omterm{def:g11:lithosphere-asthenosphere:lithosphere}{الغلاف اللدن}). وترتفع درجة الحرارة أكثر ما ترتفع بين 50
و\qty{150}{km}، بأكثر من \qty{600}{\celsius}، لأن معظم هذا
المجال يقع في \omterm{def:g11:lithosphere-asthenosphere:lithosphere}{الغلاف الصخري}؛ أما بين 150 و\qty{250}{km} فلا ترتفع
إلا بنحو \qty{45}{\celsius}.
\end{solution}

\begin{solution}{exo:g11:lithosphere-asthenosphere:7}
الجذر: $3 \times 2830/550 \approx \qty{15.4}{km}$. والسمك الإضافي:
$3 \times 3380/550 \approx \qty{18.4}{km}$، أي ارتفاع السلسلة البالغ \qty{3}{km}
مضافًا إليه جذرها.
\end{solution}

\begin{solution}{exo:g11:lithosphere-asthenosphere:8}
القشرة أسمك بمقدار $55 - 41.1 = \qty{13.9}{km}$، و$\Delta H = h
\times 3380/550$، إذن $h = 13.9 \times 550/3380 \approx \qty{2.3}{km}$.
\end{solution}

\begin{solution}{exo:g11:lithosphere-asthenosphere:9}
التناقص: $(4.470 - 4.419)/4.470 \approx \qty{1.1}{\percent}$. والقفزة:
$(4.644 - 4.419)/4.419 \approx \qty{5.1}{\percent}$. فمع العمق يزداد انضغاط
الصخر، وهذا يسرّع الموجات الزلزالية؛ والتناقص يعني
أن شيئًا آخر يغلب: فالصخر ساخن وقريب من الانصهار، وأضعف.
\end{solution}

\begin{solution}{exo:g11:lithosphere-asthenosphere:10}
(1) $(1345 - 610)/60 \approx \qty{12}{\celsius}$ لكل كيلومتر. (2) وهذا
نحو نصف التدرج عند السطح. فجزء من الحرارة يُنتج داخل
القشرة العليا؛ وعلى هذه الحرارة أيضًا أن تعبر الطبقات العليا لتخرج، وهذا يجعل
التدرج أشد انحدارًا قرب السطح. أما تحت القشرة فلا تعبر الصخر إلا الحرارة الآتية
من الأعمق، فيكون التدرج أصغر.
\end{solution}

\begin{solution}{exo:g11:lithosphere-asthenosphere:11}
عند \qty{200}{km}: $1300 + 90 = \qty{1390}{\celsius}$؛ وعند \qty{400}{km}:
$1300 + 180 = \qty{1480}{\celsius}$. فالارتفاع \qty{90}{\celsius} على امتداد
\qty{200}{km}. أما عبر \omterm{def:g11:lithosphere-asthenosphere:lithosphere}{الغلاف الصخري} فترتفع درجة الحرارة بمقدار
$1345 - 15 = \qty{1330}{\celsius}$ على امتداد \qty{100}{km} فقط: أي أكثر بنحو ثلاثين
مرة لكل كيلومتر.
\end{solution}

\begin{solution}{exo:g11:lithosphere-asthenosphere:12}
(1) $118 + 12.3\,z = 1300 + 0.45\,z$ يعطي $z = 1182/11.85 \approx
\qty{100}{km}$. (2) $118 + 21.5\,z = 1300 + 0.45\,z$ يعطي $z = 1182/21.05
\approx \qty{56}{km}$. (3) للغرب الأسخن \omterm{def:g11:lithosphere-asthenosphere:lithosphere}{غلاف صخري} أرق بكثير،
كما قيس (نحو \qty{40}{km} في أرق مواضعه)؛ أما \omterm{def:g11:lithosphere-asthenosphere:lithosphere}{الغلاف الصخري} الشديد السمك
في وسط القارة فيحتاج إلى تدرج أصغر أيضًا.
\end{solution}

\begin{solution}{exo:g11:lithosphere-asthenosphere:13}
(1) بحسب المبرهنة، $\Delta H = h\,\rho_m/(\rho_m - \rho_c)$، إذن $h =
\Delta H\,(\rho_m - \rho_c)/\rho_m$. (2) $550/3380 \approx 0.16$: فنحو
سدس أي قشرة إضافية يظهر ارتفاعًا، وخمسة أسداسها جذرًا. (3) وللجبل
الجليدي، الكسر الواقع فوق الماء $(1025 - 910)/1025 \approx 0.11$:
إنها القاعدة نفسها، (كثافة المائع $-$ كثافة الجسم الطافي)/كثافة
المائع.
\end{solution}

\begin{solution}{exo:g11:lithosphere-asthenosphere:14}
(1) $1489 - 1345 = \qty{144}{\celsius}$. (2) يتجمع قليل من المصهور في
أغشية رقيقة بين حبيبات الصخر. فلا تعود الحبيبات متشابكة
بإحكام: فيقاوم الصخر القص أقل، فتتباطأ \omterm{def:g11:seismic-structure:body-waves}{الموجات S}، ويمكنه
أن يتشوه وينساب بسهولة أكبر بكثير.
\end{solution}

\begin{solution}{exo:g11:lithosphere-asthenosphere:15}
(1) لا. فتسوية اليابسة تجعل القشرة القارية أرق قليلًا،
لكنها تبقى سميكة وخفيفة، وتبقى طافية أعلى بكثير من
القشرة الرقيقة الكثيفة لقاع المحيط؛ وتبقى المحيطات مالئة
الأحواض بين القارات. (2) هناك نوعان من القشرة، يختلفان في
السمك والكثافة: قشرة قارية سميكة خفيفة وقشرة محيطية رقيقة
كثيفة.
\end{solution}

\begin{solution}{pb:g11:lithosphere-asthenosphere:1}
\textbf{1.} $5.01 - 0.67 = \qty{4.34}{km}$.

\textbf{2.} $(3380 - 2830)/3380 \approx \qty{16}{\percent}$.

\textbf{3.} يقع تحت السطح الكسر $2830/3380 \approx 0.84$؛ أما الجبل
الجليدي فيغوص أعمق قليلًا، $910/1025 \approx 0.89$، لأن كثافة الجليد
أقرب إلى كثافة الماء من قرب كثافة القشرة إلى كثافة الوشاح.

\textbf{4.} على مدى ثواني الموجة الزلزالية، أو زمن
الزلزال، يكون الوشاح جسمًا صلبًا؛ ولا ينساب \omterm{def:g11:lithosphere-asthenosphere:lithosphere}{الغلاف اللدن} ويدع القشرة تطفو
إلا على مدى أزمنة طويلة جدًا.

\textbf{5.} تحت أعمق \omterm{def:g11:seismic-structure:moho}{موهو}، يتكوّن العمودان من صخر الوشاح نفسه
عند الأعماق نفسها، فيتساوى وزنهما هناك؛ وكل الفروق
تقع فوق ذلك العمق ويجب أن تُحتسب.

\textbf{6.} القارة المتوسطة: قشرة من \qty{0.67}{km} فوق مستوى سطح البحر إلى
\qty{40.4}{km} تحته، ثم وشاح حتى \qty{70}{km}. والهضبة:
قشرة من \qty{5.01}{km} فوق مستوى سطح البحر حتى \omterm{def:g11:seismic-structure:moho}{موهو} أعمق بمقدار
الجذر، ثم وشاح حتى \qty{70}{km}.

\textbf{7.} يحوي عمود الهضبة قشرةً تزيد بمقدار $h + r$، ويحوي عمود القارة
بدلًا منها وشاحًا على السمك $r$: $\rho_c\,(h + r) = \rho_m\,r$،
إذن $r = h\,\rho_c/(\rho_m - \rho_c)$ و$\Delta H = h + r =
h\,\rho_m/(\rho_m - \rho_c)$.

\textbf{8.} $3380/550 \approx 6.15$.

\textbf{9.} $\Delta H = 4.34 \times 6.15 \approx \qty{26.7}{km}$؛ والجذر
$26.7 - 4.34 \approx \qty{22.3}{km}$.

\textbf{10.} يقع \omterm{def:g11:seismic-structure:moho}{موهو} القارة المتوسطة على عمق $41.1 - 0.67 \approx
\qty{40.4}{km}$ تحت مستوى سطح البحر؛ \omterm{def:g11:seismic-structure:moho}{وموهو} الهضبة على عمق $40.4 + 22.3 \approx
\qty{62.8}{km}$ تحت مستوى سطح البحر.

\textbf{11.} $41.1 + 26.7 \approx \qty{67.8}{km}$.

\textbf{12.} كتلة كيلومتر واحد من الصخر فوق متر مربع هي
$\rho \times \qty{1000}{m^3}$: أي \num{2830} طنًا من القشرة أو \num{3380}
طنًا من الوشاح. القارة: $41.1 \times 2830 + (70 - 40.43) \times 3380
\approx \num{116300} + \num{99900} \approx \num{216300}$ طن. والهضبة:
$67.77 \times 2830 + (70 - 62.76) \times 3380 \approx \num{191800} +
\num{24500} \approx \num{216300}$ طن. فالكتلتان متساويتان.

\textbf{13.} $41.1 - 0.67 \times 6.15 \approx \qty{37.0}{km}$.

\textbf{14.} $67.8/41.1 \approx 1.65$ مرة.

\textbf{15.} $3300/470 \approx 7.02$، و$\Delta H \approx 4.34 \times 7.02
\approx \qty{30.5}{km}$، أي قشرة سمكها نحو \qty{71.6}{km}. فتغيّر قدره
نحو \qty{2}{\percent} في كثافة الوشاح يغيّر السمك بنحو
\qty{4}{km}: فالنتيجة تتوقف على الكثافتين توقفًا ملحوظًا، لكنها تبقى
قريبة من \qty{70}{km}.

\textbf{16.} يعطي النموذج نحو \qty{68}{km}، والموجات الزلزالية نحو
\qty{70}{km}: فهما متفقان في حدود بضعة في المئة.

\textbf{17.} $h = 3.23 - 0.67 = \qty{2.56}{km}$، و$\Delta H = 2.56 \times
6.15 \approx \qty{15.7}{km}$، أي قشرة سمكها نحو \qty{56.8}{km}، تزيد بنحو
\qty{7}{km} على المقيس.

\textbf{18.} أرق. فالعمود الذي جزؤه العلوي أخف من المتوسط
يزن أقل للسمك نفسه، فيطفو أعلى: والقشرة الأخف
تحتاج إلى جذر أصغر لترتفع إلى ارتفاع معين. فيمكن إذن أن تكون القشرة الحقيقية
أرق من قشرة النموذج، كما قيس.

\textbf{19.} على مدى أزمنة طويلة ينساب \omterm{def:g11:lithosphere-asthenosphere:lithosphere}{الغلاف اللدن}: فهو يتنحى تحت
عمود ثقيل ويعود تحت عمود خفيف، حتى تتوازن الأعمدة.

\textbf{20.} \textbf{يبلغ سمك القشرة تحت التبت نحو \qty{70}{km}}، أي أكثر
بكثير من سمك قشرة القارة المتوسطة البالغ \qty{41}{km}: فالهضبة ترتفع خمسة
كيلومترات لأن قشرتها السميكة الخفيفة تطفو على الوشاح الأكثف،
بجذر يزيد على \qty{20}{km}.
\end{solution}
